\scanpage{102}
Ce qui signifie aussi que $\forall$ \inserted{(ou un)} scindage de 31 --- p.ex. défini par un pt de $U_\eta$ rat./$\unsure{K|k}$ --- induit un hom $T\longrightarrow\pi$)

\noindent c) L'action de $T$ sur $\pi_{\mathrm{ab}}=H_1(U_{\bar\eta})$ est triviale.

\noindent d) Idem pour $H_1(U_{\bar\eta},\unsure{T_\ell|\mathbf Z_\ell})$, \unsure{si|où} $\ell$ est premier à la car. résiduelle. (\uline{NB} \deleted{\unsure{Pour car.} $p>0$} \inserted{$\ell=p>0$}, on doit pouvoir l'exprimer par un critère sur les groupes $p$-divisibles\ldots)

\noindent e) (En termes d'une section de (30), définissant une extension de $\mathcal E_k$ par $\pi$, et supposant maintenant $k$ de type fini sur $\mathbf Q$) Ladite extension est admissible (quitte à passer à un sous-groupe ouvert, \inserted{de $\mathcal E_k$,} au besoin), i.e. peut être définie par une courbe alg. sur \inserted{$k$, ou sur} une extension finie de $k$.

\noindent f) (En termes d'une section de (30) et d'une de (31), permettant de faire opérer $\mathcal E_k$ sur $\pi$) : Cette action \inserted{n'}est \inserted{pas trop} loufoque (à définir convenablement !)
\margin{101}
